\chapter{Planificación temporal}
\label{ch:definicion-del-problema}
\paragraph*{}La planificación temporal de este Trabajo Fin de Grado se organiza en varias fases, pensadas para abordar de forma ordenada la ampliación y mejora de la aplicación existente. La verdad es que, en un proyecto de este tipo, contar con una planificación clara ayuda mucho: permite avanzar sin perder de vista los objetivos principales y evita que las mejoras técnicas queden desconectadas entre sí. Cada fase se centra en un conjunto de tareas concretas, desde la preparación inicial del entorno hasta la validación final de los resultados obtenidos.

\begin{enumerate}
    \item \textbf{Preparación del entorno y revisión del sistema existente} \\
    En esta primera fase se realizará una revisión completa del TFG original y de la aplicación desarrollada previamente. El objetivo será comprender con calma su funcionamiento interno, sus dependencias, la estructura del código y la forma en la que se organizan los datos. Además, se configurará el entorno de trabajo en el equipo local, incluyendo las herramientas de desarrollo necesarias, las bibliotecas de Python, el entorno de Dash y los archivos de datos empleados por la aplicación. \\
    Esta etapa servirá como punto de partida técnico. Antes de incorporar nuevas funcionalidades, resulta necesario asegurarse de que el sistema base funciona correctamente y de que se conocen sus límites.
    
    \item \textbf{Análisis inicial y planificación} \\
    En esta fase se identificarán las limitaciones del sistema original y las áreas que requieren una mejora más clara. También se definirán los nuevos requisitos funcionales, no funcionales y de información, así como los objetivos técnicos y metodológicos del proyecto. Además, se elaborará un plan de trabajo que permita repartir mejor los recursos, los tiempos y las herramientas que se utilizarán durante el desarrollo.
    
    \item \textbf{Diseño de mejoras y ampliaciones} \\
    A partir del análisis previo, se planificarán las mejoras principales de la aplicación. En esta parte se estudiará cómo incorporar el soporte multilingüe, la implementación de gráficos adicionales basados en el Test de Log-Rank y en la regresión de Cox, y los cambios necesarios en la interfaz para que la experiencia de uso resulte más clara. También se valorará el tratamiento de atributos no binarios del dataset, junto con la integración de una inteligencia artificial más rápida y eficiente.
    
    \item \textbf{Desarrollo e implementación} \\
    Durante esta fase se llevará a cabo la programación de las mejoras diseñadas. Será, previsiblemente, la parte más extensa del trabajo, ya que aquí las ideas definidas en las fases anteriores se convierten en funcionalidades reales. Se implementarán nuevos módulos, se integrarán funcionalidades interactivas y se añadirán modelos de análisis de supervivencia adicionales. Asimismo, se actualizará la interfaz de usuario en \underline{Dash} y se optimizará el código para mejorar el rendimiento general del sistema.
    
    \item \textbf{Pruebas y validación} \\
    En esta fase se realizarán pruebas para comprobar la estabilidad y fiabilidad del sistema mejorado. No basta con que la aplicación funcione una vez; debe responder de forma coherente ante distintos escenarios. Por ello, se evaluará el comportamiento de las nuevas funcionalidades, la precisión de los modelos estadísticos, la exportación de informes e imágenes y la velocidad de respuesta de la inteligencia artificial. Se emplearán datos reales y, cuando sea necesario, datos de prueba para comprobar la calidad de los resultados.
    
    \item \textbf{Documentación técnica y manuales} \\
    Se elaborará la documentación completa del proyecto, incluyendo el manual técnico, el manual de usuario y una guía de instalación y ejecución. En estos documentos se explicará la estructura de la aplicación, el funcionamiento de las nuevas funcionalidades y las instrucciones necesarias para utilizar el sistema. Además, se documentará el código fuente de forma clara, de manera que el trabajo pueda entenderse y mantenerse en el futuro.
    
    \item \textbf{Evaluación final y conclusiones} \\
    En la última fase se analizará el cumplimiento de los objetivos planteados y se evaluará el impacto de las mejoras implementadas sobre la aplicación original. También se expondrán las conclusiones generales del trabajo, los problemas encontrados durante el desarrollo y las posibles líneas de mejora que podrían abordarse en el futuro.
\end{enumerate}

\newpage

\textbf{Cronograma de trabajo:}

A continuación se presenta una estimación de la distribución temporal del proyecto, con una duración total aproximada de 300 horas:


\begin{center}
\rowcolors{2}{blue!5!white}{white}
\begin{tabular}{|p{9cm}|c|}
\hline
\rowcolor{blue!20!white}
\textbf{FASE} & \textbf{TIEMPO (HORAS)} \\
\hline
\textcolor{blue!60!black}{\textbf{Preparación}} & 30 \\
\hline
\textcolor{blue!60!black}{\textbf{Análisis y diseño}} & 60 \\
\hline
\textcolor{blue!60!black}{\textbf{Desarrollo e implementación}} & 120 \\
\hline
\textcolor{blue!60!black}{\textbf{Pruebas y validación}} & 50 \\
\hline
\textcolor{blue!60!black}{\textbf{Documentación y manuales}} & 25 \\
\hline
\textcolor{blue!60!black}{\textbf{Evaluación final y conclusiones}} & 15 \\
\hline
\rowcolor{blue!20!white}
\textbf{TOTAL ESTIMADO} & \textbf{300 HORAS} \\
\hline
\end{tabular}
\end{center}

\paragraph*{}La tabla muestra cómo se reparten las horas de trabajo entre las distintas fases del proyecto. Como es lógico, la parte que más dedicación requiere es la de desarrollo e implementación, ya que en ella se concentran la mayoría de las mejoras de la aplicación. También tienen un peso importante las fases de análisis, diseño y validación, porque ayudan a definir bien el camino antes de programar y a comprobar después que los resultados obtenidos son fiables. En conjunto, esta distribución ofrece una visión realista del esfuerzo previsto y permite organizar el proyecto con mayor claridad.
